% ============================================================================
%  ANEXO A: CÓDIGO Y FIRMAS DE SEGURIDAD DEL LABORATORIO
% ============================================================================

\section{Anexo A: Archivo de Firmas de Seguridad (homelab.rules)}

A continuación se transcribe el conjunto de reglas personalizadas implementadas en el servidor \texttt{dc1} (\texttt{/etc/suricata/rules/homelab.rules}) para la detección de anomalías y bloqueo activo de intrusiones:

\begin{lstlisting}[language=bash]
# ===================================================================
# homelab.rules — Reglas personalizadas de IDS/IPS para dc1 (Sudoers LAN)
# ===================================================================

# 1. PRUEBAS Y VALIDACION DE LABORATORIO (IDS vs IPS)
drop tcp any any -> $HOME_NET any (msg:"[HOMELAB-IPS-DROP] Bloqueo en tiempo real: Firma de ataque detectada (TEST_EXPLOIT_PAYLOAD_DROP)"; flow:to_server,established; content:"TEST_EXPLOIT_PAYLOAD_DROP"; nocase; classtype:attempted-admin; sid:9000001; rev:1;)
drop tcp any any -> $HOME_NET $HTTP_PORTS (msg:"[HOMELAB-IPS-DROP] Bloqueo HTTP: Carga maliciosa en aplicacion web (MALICIOUS_SIMULATION_ATTACK)"; flow:to_server,established; content:"MALICIOUS_SIMULATION_ATTACK"; nocase; classtype:web-application-attack; sid:9000002; rev:1;)

alert tcp any any -> $HOME_NET any (msg:"[HOMELAB-IDS-ALERT] Sonda de prueba detectada (TEST_PROBE_ALERT_ONLY)"; flow:to_server,established; content:"TEST_PROBE_ALERT_ONLY"; nocase; classtype:misc-activity; sid:9000010; rev:1;)
alert icmp any any -> $HOME_NET any (msg:"[HOMELAB-IDS-ALERT] Ping ICMP Echo Request recibido en dc1"; itype:8; classtype:misc-activity; sid:9000011; rev:1;)

# 2. DETECCION DE ESCANEOS Y RECONOCIMIENTO DE RED (NIDS)
alert tcp any any -> $HOME_NET any (msg:"[HOMELAB-IDS-ALERT] Reconocimiento de red: Escaneo TCP NULL detectado"; flags:0; classtype:attempted-recon; sid:9000020; rev:1;)
alert tcp any any -> $HOME_NET any (msg:"[HOMELAB-IDS-ALERT] Reconocimiento de red: Escaneo TCP XMAS detectado"; flags:FPU; classtype:attempted-recon; sid:9000021; rev:1;)
alert tcp any any -> $HOME_NET any (msg:"[HOMELAB-IDS-ALERT] Posible escaneo de puertos SYN Scan en progreso"; flags:S,12; threshold:type both, track by_src, count 25, seconds 3; classtype:attempted-recon; sid:9000022; rev:1;)

# 3. PROTECCION DE SERVICIOS CRITICOS (Active Directory / Samba DC)
alert tcp any any -> $HOME_NET 88 (msg:"[HOMELAB-IDS-ALERT] Sondeo continuo o sospechoso al servicio Kerberos del AD DC"; flags:S; threshold:type both, track by_src, count 15, seconds 5; classtype:attempted-recon; sid:9000030; rev:1;)
alert tcp any any -> $HOME_NET 389 (msg:"[HOMELAB-IDS-ALERT] Multiples intentos de conexion a LDAP (posible enumeracion de cuentas AD)"; flags:S; threshold:type both, track by_src, count 20, seconds 5; classtype:attempted-recon; sid:9000031; rev:1;)
alert tcp any any -> $HOME_NET 445 (msg:"[HOMELAB-IDS-ALERT] Trafico SMB repetitivo detectado hacia dc1 (posible sondeo de shares)"; flags:S; threshold:type both, track by_src, count 20, seconds 10; classtype:suspicious-login; sid:9000032; rev:1;)

# 4. ACCESO Y FUERZA BRUTA SSH (Puerto 22)
alert tcp any any -> $HOME_NET 22 (msg:"[HOMELAB-IDS-ALERT] Intento de conexion SSH entrante a dc1"; flags:S; classtype:suspicious-login; sid:9000040; rev:1;)
alert tcp any any -> $HOME_NET 22 (msg:"[HOMELAB-IDS-ALERT] ALERTA CRITICA: Posible ataque de fuerza bruta SSH contra dc1"; flags:S; threshold:type both, track by_src, count 6, seconds 30; classtype:attempted-admin; sid:9000041; rev:1;)

# 5. PROTECCION DE APLICACIONES WEB (Nginx / Intranet / Portal)
alert tcp any any -> $HOME_NET $HTTP_PORTS (msg:"[HOMELAB-IDS-ALERT] Ataque Web: Patron de Inyeccion SQL (UNION SELECT)"; flow:to_server,established; content:"UNION"; nocase; content:"SELECT"; nocase; distance:0; classtype:web-application-attack; sid:9000050; rev:1;)
alert tcp any any -> $HOME_NET $HTTP_PORTS (msg:"[HOMELAB-IDS-ALERT] Ataque Web: Intento de Path Traversal hacia /etc/passwd"; flow:to_server,established; content:"/etc/passwd"; nocase; classtype:web-application-attack; sid:9000051; rev:1;)
alert tcp any any -> $HOME_NET $HTTP_PORTS (msg:"[HOMELAB-IDS-ALERT] Ataque Web: Intento de explotacion Log4Shell (jndi:ldap)"; flow:to_server,established; content:"${jndi:"; nocase; classtype:web-application-attack; sid:9000052; rev:1;)
\end{lstlisting}

\newpage
\section{Anexo B: Script Alternador entre Modos IDS e IPS (ips-toggle.sh)}

A continuación se presenta el código de la utilidad \texttt{server/suricata/ips-toggle.sh} desarrollada para alternar de forma transparente el acoplamiento con \texttt{nftables}:

\begin{lstlisting}[language=bash]
#!/bin/bash
# ips-toggle.sh — Alternador entre Modo IDS y Modo IPS
set -euo pipefail
MODE="${1:-}"

has_nft_queue() {
    nft list chain inet filter input 2>/dev/null | grep -q "queue num 0"
}

remove_nft_queue() {
    if has_nft_queue; then
        local handles
        handles="$(nft -a list chain inet filter input 2>/dev/null | grep "queue num 0" | awk '{print $NF}')"
        for h in $handles; do
            nft delete rule inet filter input handle "$h" 2>/dev/null || true
        done
    fi
}

add_nft_queue() {
    if ! has_nft_queue; then
        nft insert rule inet filter input queue num 0 bypass
    fi
}

case "$MODE" in
    --ids)
        remove_nft_queue
        sed -i 's|SURICATA_OPTS=.*|SURICATA_OPTS="-i eno1"|' /etc/default/suricata
        systemctl restart suricata
        echo "Modo IDS activo (AF_PACKET pasivo en eno1)"
        ;;
    --ips)
        add_nft_queue
        sed -i 's|SURICATA_OPTS=.*|SURICATA_OPTS="-q 0"|' /etc/default/suricata
        systemctl restart suricata
        echo "Modo IPS activo (NFQUEUE 0 inline en nftables)"
        ;;
    --status)
        echo "Estado de Suricata:"
        systemctl is-active suricata
        ;;
    *)
        echo "Uso: sudo bash $0 {--ids | --ips | --status}"
        exit 1
        ;;
esac
\end{lstlisting}
